% Created 2026-03-05 jue 10:52
% Intended LaTeX compiler: pdflatex
\documentclass[a4paper]{article}
\usepackage[utf8]{inputenc}
\usepackage[T1]{fontenc}
\usepackage{graphicx}
\usepackage{longtable}
\usepackage{wrapfig}
\usepackage{rotating}
\usepackage[normalem]{ulem}
\usepackage{amsmath}
\usepackage{amssymb}
\usepackage{capt-of}
\usepackage{hyperref}
\usepackage{color}
\usepackage{listings}
\usepackage[spanish]{babel}
\usepackage[usenames,dvipsnames]{color} % Required for custom colors
\renewcommand{\ttdefault}{pcr} % MONOESPACIO CON NEGRIT
\usepackage{lastpage}
\usepackage{listings}
\usepackage{listingsutf8}
\renewcommand{\lstlistingname}{Listado}
\lstset{frame=single,inputencoding=utf8,basicstyle=\scriptsize\ttfamily,showstringspaces=false,numbers=none}
\definecolor{MyDarkGreen}{rgb}{0.0,0.4,0.0} % This is the color used for comments
\lstset{ breaklines=true, postbreak=\mbox{\textcolor{red}{$\hookrightarrow$}\space}, keywordstyle=\bfseries, keywordstyle=[1]\color{Blue}\bfseries,  keywordstyle=[2]\color{Purple}\bfseries,  keywordstyle=[3]\color{Blue}\underbar,   identifierstyle=,   commentstyle=\usefont{T1}{pcr}{m}{sl}\color{MyDarkGreen}\small,   stringstyle=\color{Purple},   showstringspaces=false,   tabsize=2,   morecomment=[l][\color{Blue}]{...} }
\lstset{literate=  {á}{{\'a}}1 {é}{{\'e}}1 {í}{{\'i}}1 {ó}{{\'o}}1 {ú}{{\'u}}1   {Á}{{\'A}}1 {É}{{\'E}}1 {Í}{{\'I}}1 {Ó}{{\'O}}1 {Ú}{{\'U}}1   {à}{{\`a}}1 {è}{{\`e}}1 {ì}{{\`i}}1 {ò}{{\`o}}1 {ù}{{\`u}}1   {À}{{\`A}}1 {È}{{\'E}}1 {Ì}{{\`I}}1 {Ò}{{\`O}}1 {Ù}{{\`U}}1   {ä}{{\"a}}1 {ë}{{\"e}}1 {ï}{{\"i}}1 {ö}{{\"o}}1 {ü}{{\"u}}1   {Ä}{{\"A}}1 {Ë}{{\"E}}1 {Ï}{{\"I}}1 {Ö}{{\"O}}1 {Ü}{{\"U}}1   {â}{{\^a}}1 {ê}{{\^e}}1 {î}{{\^i}}1 {ô}{{\^o}}1 {û}{{\^u}}1   {Â}{{\^A}}1 {Ê}{{\^E}}1 {Î}{{\^I}}1 {Ô}{{\^O}}1 {Û}{{\^U}}1   {œ}{{\oe}}1 {Œ}{{\OE}}1 {æ}{{\ae}}1 {Æ}{{\AE}}1 {ß}{{\ss}}1   {ű}{{\H{u}}}1 {Ű}{{\H{U}}}1 {ő}{{\H{o}}}1 {Ő}{{\H{O}}}1   {ç}{{\c c}}1 {Ç}{{\c C}}1 {ø}{{\o}}1 {å}{{\r a}}1 {Å}{{\r A}}1   {€}{{\euro}}1 {£}{{\pounds}}1 {«}{{\guillemotleft}}1   {»}{{\guillemotright}}1 {ñ}{{\~n}}1 {Ñ}{{\~N}}1 {¿}{{?`}}1 }
\usepackage{caption}
\usepackage{attachfile2}
\usepackage[margin=1.5cm,includeheadfoot,includehead,includefoot]{geometry}
\hypersetup{colorlinks,linkcolor=black}
\usepackage{fancyhdr}
\pagestyle{fancyplain}
\chead{}
\lhead{}
\rhead{}
\cfoot{}
\lfoot{\begin{footnotesize}alvaro.gonzalezsotillo@educa.madrid.org\end{footnotesize}}
\rfoot{\begin{footnotesize}\thepage / \pageref{LastPage}\end{footnotesize}}
\usepackage[skins]{tcolorbox}
\usepackage{multicol}
\usepackage{changepage} %ajdustwidth
\usepackage{fancybox}
\usepackage{attachfile2}
\lhead{Extraordinaria 2021 (es el \\lhead)}
\rhead{Administración y Gestión de Bases de Datos (es el \\rhead)}
\lhead{Clean code}
\rhead{Entornos de desarrollo}
\usepackage{svg}
\usepackage{letltxmacro}
\LetLtxMacro{\originalincludegraphics}{\includegraphics}
\renewcommand{\includegraphics}[2][]{\IfFileExists{#2.pdf}{\originalincludegraphics[#1]{#2.pdf}}{\originalincludegraphics[#1]{#2}}}
\LetLtxMacro{\originalincludesvg}{\includesvg}
\renewcommand{\includesvg}[2][]{\IfFileExists{#2.pdf}{\originalincludegraphics[#1]{#2.pdf}}{\originalincludegraphics[#1]{#2.svg.pdf}}}
\usepackage{comment}
\excludecomment{NOTES}
\date{}
\title{}
\hypersetup{
 pdfauthor={Álvaro González Sotillo},
 pdftitle={},
 pdfkeywords={},
 pdfsubject={},
 pdfcreator={Emacs 30.2 (Org mode 9.7.11)}, 
 pdflang={Spanish}}
\begin{document}

\captionsetup{font=scriptsize}

\textattachfile{/datos-luks/datos-alvaro/apuntes-clase/entornos-de-desarrollo/apuntes/4/ED-04-practica-cleancode.org}{} \vspace{-1em}


\setlength{\parindent}{0em}
\setlength{\parskip}{1em}

\newtcolorbox{Aviso}[1][Aviso]{
  enhanced,
  colback=gray!5!white,
  colframe=gray!75!black,fonttitle=\bfseries,
  colbacktitle=gray!85!black,
  attach boxed title to top left={yshift=-2mm,xshift=2mm},
  title=#1
}

\newtcolorbox{cuadrito}[1][Ignorado]{
  %drop shadow southeast,
  enhanced jigsaw,
  colback=white,
}


\newcommand{\StudentData}{
  \begin{cuadrito}[1\textwidth]
    \vspace{0.3cm}
    \large{
      \textbf{Apellidos:} \hrulefill \\
      \textbf{Nombre:} \hrulefill \\
      \textbf{Fecha:} \hrulefill \hspace{1cm} \textbf{Usuario:} \hrulefill \\
    }
    \vspace{-0.2cm}
  \end{cuadrito}
}

\newcommand{\StudentDataDniFirma}{
  \begin{cuadrito}[1\textwidth]
    \vspace{0.3cm}
    \large{
      \textbf{Apellidos:} \hrulefill \\
      \textbf{Nombre:} \hrulefill  \\
      \textbf{Dni:} \hrulefill \hspace{8cm} \\
      \\
      \textbf{Firma:} \hrulefill \hspace{8cm} \\
    }
    \vspace{-0.2cm}
  \end{cuadrito}
}
\section{Objetivo de la práctica}
\label{sec:org0000000}
En esta práctica el alumno seguirá trabajando indirectamente el control de versiones, además de:
\begin{itemize}
\item Uso básico de \emph{javadoc}
\item Refactorización usando un IDE
\item Seguimiento de guías de estilo y eliminación de \emph{bad smells}
\end{itemize}

Esta práctica será tenida en cuenta para el RA4.
\section{Enunciado}
\label{sec:org000000f}
Este programa resuelve el problema del \href{http://es.wikipedia.org/wiki/Pentomin\%C3\%B3}{pentominó}. El código fuente está bastante bien, pero es necesario mejorarlo en algunos aspectos. En concreto, hay que mejorar ciertas partes internas del código y no está documentado adecuadamente. El código fuente inicial está en \url{https://codeberg.org/alvarogonzalezsotillo/pentominos}
\subsection{Documentación (2 puntos)}
\label{sec:org0000003}
Hay clases, atributos y métodos sin documentar. Se debe completar la documentación (en español o en inglés, indistintamente)

\begin{itemize}
\item Todas las clases tendrán \texttt{@author} (usa tu nombre y apellido). Existen desde la versión 1.0.
\item Todas las clases tendrán alguna referencia \texttt{@see} a otra que se use o por la que sea usada.
\item En la clase \texttt{Pentominos}  
\begin{itemize}
\item Todos los métodos tendrán \texttt{@param} para todos sus atributos. Todos los métodos con retorno tendrán \texttt{@return}
\item Cada método tendrá una referencia \texttt{@link} a algún método que sea llamado dentro de él. Si no llama a ningún método, tendrá una referencia a un método que lo llame.
\end{itemize}
\end{itemize}
\subsection{Refactorización (4 puntos)}
\label{sec:org0000006}
Realiza estos cambios con las herramientas de refactorización de Intellij:
\begin{itemize}
\item Las normas de codificación de nuestra metodología indican que los atributos de las clases comienzan con el prefijo \texttt{m\_}. Todos los atributos de la clase \texttt{MosaicPanel} se modificarán.
\item El \texttt{PentominosPanel.menuHandler} es demasiado largo. Cada acción se realizará en un método aparte, en vez de dentro del \texttt{if} encadenado.
\item El método \texttt{PentominosPanel.doSaveImage} se dividirá en dos:
\begin{itemize}
\item Una parte pedirá un fichero y comprobará que no existe
\item La otra parte utilizará ese fichero para grabar la imagen.
\end{itemize}
\item Los atributos \texttt{rows} y \texttt{columns} de \texttt{MosaicPanel} no se accederán directamente. Se utilizarán métodos \texttt{set} y \texttt{get} para ello.
\item Las clases deberán formar parte del paquete \texttt{edu.hws.godel}.
\end{itemize}
\subsection{Clean code (3)}
\label{sec:org0000009}
\begin{itemize}
\item Consigue que la clase MakeIcons no tenga ningún aviso de Sonarqube.
\item Puedes \href{https://stackoverflow.com/questions/10971968/turning-sonar-off-for-certain-code}{ignorar el aviso} \texttt{S106}
\end{itemize}
\subsection{Javadoc (1 punto)}
\label{sec:org000000c}
\begin{itemize}
\item Genera la documentación \emph{javadoc} del proyecto
\item Se dejará en el directorio \texttt{javadoc}
\end{itemize}
\section{Normas de entrega}
\label{sec:org0000012}
\begin{enumerate}
\item Se creará un repositorio en Codeberg, siendo inicialmente un \emph{fork} de \url{https://codeberg.org/alvarogonzalezsotillo/pentominos}
\item La práctica se realizará en la rama \texttt{2026-refactorizacion-documentacion}
\item Se modificará el código sobre ese repositorio hasta ajustarse a lo pedido en la práctica.
\item Se creará un \emph{tag} con la versión para entregar, de nombre \texttt{2026-entrega}
\item En el aula virtual, se subirá:
\begin{itemize}
\item La URL del repositorio
\item El identificador del \emph{commit} del \emph{tag} de la entrega.
\end{itemize}
\item Si la práctica se realiza por parejas, todos los integrantes subirán el resultado al aula virtual, y en la historia del repositorio podrá verse que el reparto de trabajo ha sido equilibrado
\end{enumerate}
\section{Qué se valorará}
\label{sec:org0000015}
\begin{itemize}
\item Que haya un \emph{commit} por cada subapartado.
\item Que el código final se ajuste a lo pedido
\item Que el código funcione (clases Main: \texttt{Pentominos} y \texttt{MakeIcons})
\end{itemize}
\end{document}
